\documentclass[12pt,a4paper]{article}

\usepackage[a4paper,margin=1in]{geometry}
\usepackage{graphicx}
\usepackage{float}
\usepackage{booktabs}
\usepackage{array}
\usepackage{hyperref}
\usepackage{titlesec}
\usepackage{setspace}
\usepackage{fancyhdr}
\usepackage{xcolor}
\usepackage{tikz}
\usetikzlibrary{calc}

\setstretch{1.25}

\pagestyle{fancy}
\fancyhf{}
\lhead{Embedded Systems}
\rhead{Calculator Project}
\cfoot{\thepage}

\begin{document}

%%%%%%%%%%%%%%%%%%%%%%%%%%%%%%%%%%%%%%%%%%%%%%%%%%
% TITLE PAGE
%%%%%%%%%%%%%%%%%%%%%%%%%%%%%%%%%%%%%%%%%%%%%%%%%%

\begin{titlepage}
\begin{tikzpicture}[remember picture,overlay]
\draw[line width=2pt]
    ($(current page.north west)+(1cm,-1cm)$)
    rectangle
    ($(current page.south east)+(-1cm,1cm)$);
\end{tikzpicture}

\centering

\vspace*{2cm}

{\Huge\bfseries SCIENTIFIC CALCULATOR}

\vspace{0.3cm}

{\Large Arduino Uno Based Breadboard Implementation}

\vspace{1.5cm}

{\Large Project Report}

\vspace{1.8cm}

\begin{tabular}{ll}
\textbf{Microcontroller} & Arduino Uno\\
\textbf{Display} & 16×2 LCD Display\\
\textbf{Input} & 5×5 Push Button Matrix\\
\textbf{Programming} & Arduino IDE\\
\textbf{Implementation} & Breadboard Prototype\\
\end{tabular}

\vfill

\textbf{\Large Submitted By}

\vspace{0.3cm}

Kalpana Kumari

June, 2026

\end{titlepage}

%%%%%%%%%%%%%%%%%%%%%%%%%%%%%%%%%%%%%%%%%%%%%%%%%%
% CONTENTS
%%%%%%%%%%%%%%%%%%%%%%%%%%%%%%%%%%%%%%%%%%%%%%%%%%

\tableofcontents

\newpage

%%%%%%%%%%%%%%%%%%%%%%%%%%%%%%%%%%%%%%%%%%%%%%%%%%
% INTRODUCTION
%%%%%%%%%%%%%%%%%%%%%%%%%%%%%%%%%%%%%%%%%%%%%%%%%%

\section{Introduction}

A scientific calculator is an embedded electronic system capable of performing arithmetic and mathematical operations using a programmable microcontroller. Unlike a basic calculator, it supports additional mathematical functions while providing a simple user interface through a keypad and LCD display.

In this project, a scientific calculator was implemented using an Arduino Uno, a 16×2 LCD display, and a 5×5 push-button matrix assembled entirely on a breadboard. The calculator accepts user input through the keypad, processes the entered expression using the Arduino microcontroller, and displays the result on the LCD.

Since the objective of this project was hardware implementation and testing, the complete circuit was constructed and verified on a breadboard without transferring it to a permanent PCB.

%%%%%%%%%%%%%%%%%%%%%%%%%%%%%%%%%%%%%%%%%%%%%%%%%%
% OBJECTIVES
%%%%%%%%%%%%%%%%%%%%%%%%%%%%%%%%%%%%%%%%%%%%%%%%%%

\section{Objectives}

The objectives of this project are:

\begin{itemize}

\item To design a scientific calculator using Arduino Uno.

\item To interface a 16×2 LCD display.

\item To implement a 5×5 button matrix for user input.

\item To perform basic arithmetic and scientific calculations.

\item To understand keypad scanning techniques.

\item To verify the complete circuit using a breadboard.

\end{itemize}

%%%%%%%%%%%%%%%%%%%%%%%%%%%%%%%%%%%%%%%%%%%%%%%%%%
% COMPONENTS
%%%%%%%%%%%%%%%%%%%%%%%%%%%%%%%%%%%%%%%%%%%%%%%%%%

\section{Components Required}

\begin{table}[H]

\centering

\caption{Components Used}

\begin{tabular}{|p{8cm}|c|}

\hline

\textbf{Component} & \textbf{Quantity}\\

\hline

Arduino Uno & 1\\

\hline

16×2 LCD Display & 1\\

\hline

Push Buttons & 25\\

\hline

Breadboard & 1\\

\hline

Potentiometer (10k$\Omega$) & 1\\

\hline

Jumper Wires & As Required\\

\hline

USB Cable & 1\\

\hline

\end{tabular}

\end{table}

%%%%%%%%%%%%%%%%%%%%%%%%%%%%%%%%%%%%%%%%%%%%%%%%%%
% WORKING PRINCIPLE
%%%%%%%%%%%%%%%%%%%%%%%%%%%%%%%%%%%%%%%%%%%%%%%%%%

\section{Working Principle}

The Arduino Uno acts as the central controller of the scientific calculator. A 5×5 matrix of push buttons is continuously scanned to detect key presses while minimizing the number of I/O pins required.

When a button is pressed, the Arduino identifies its corresponding row and column, interprets the selected operation, performs the required calculation, and displays the result on the 16×2 LCD.

The LCD serves as the user interface by displaying the entered expression as well as the computed output. Since all hardware components are assembled on a breadboard, the circuit can easily be modified and tested during development.

%%%%%%%%%%%%%%%%%%%%%%%%%%%%%%%%%%%%%%%%%%%%%%%%%%
% CIRCUIT DIAGRAM
%%%%%%%%%%%%%%%%%%%%%%%%%%%%%%%%%%%%%%%%%%%%%%%%%%

\section{Circuit Diagram}

Figure given below illustrates the complete circuit diagram of the scientific calculator. The circuit consists of an Arduino Uno interfaced with a 16×2 LCD display, a 5×5 push-button matrix for user input, and the required power connections. The Arduino continuously scans the keypad matrix, processes the selected mathematical operation, and displays the corresponding result on the LCD.

\begin{figure}[H]
    \centering
    % Replace calculator_schematic.png with your actual image filename
    \includegraphics[width=\textwidth]{calculator.png}
    \caption{Circuit Diagram of the Scientific Calculator}
    \label{fig:circuit}
\end{figure}

%%%%%%%%%%%%%%%%%%%%%%%%%%%%%%%%%%%%%%%%%%%%%%%%%%
% BREADBOARD IMPLEMENTATION
%%%%%%%%%%%%%%%%%%%%%%%%%%%%%%%%%%%%%%%%%%%%%%%%%%

\section{Breadboard Implementation}

The complete scientific calculator circuit was assembled and tested on a breadboard. This approach allowed easy modification of the circuit connections during testing and debugging.

The Arduino Uno was connected to a 16×2 LCD display, a 5×5 push-button matrix, and the necessary power connections using jumper wires. Each button press was detected by scanning the keypad matrix, after which the corresponding mathematical operation was executed and the result displayed on the LCD.

The breadboard implementation successfully demonstrated the complete functionality of the calculator without requiring a permanent PCB.

\begin{figure}[H]
    \centering
    % Replace calculator_breadboard.jpg with your image name
    \includegraphics[width=0.75\textwidth]{calculator.jpeg}
    \caption{Breadboard Implementation of the Scientific Calculator}
\end{figure}

%%%%%%%%%%%%%%%%%%%%%%%%%%%%%%%%%%%%%%%%%%%%%%%%%%
% CIRCUIT CONNECTIONS
%%%%%%%%%%%%%%%%%%%%%%%%%%%%%%%%%%%%%%%%%%%%%%%%%%

\section{Circuit Connections}

The major hardware connections used in the calculator are summarized below.

\begin{itemize}

\item Arduino Uno serves as the main processing unit.

\item A 16×2 LCD is interfaced in 4-bit mode to display inputs and results.

\item A 5×5 matrix keypad is connected to the Arduino using row and column scanning.

\item A 10 k$\Omega$ potentiometer is used to adjust the LCD contrast.

\item USB power supply provides regulated 5 V to the complete circuit.

\item Jumper wires are used to establish all electrical connections on the breadboard.

\end{itemize}

%%%%%%%%%%%%%%%%%%%%%%%%%%%%%%%%%%%%%%%%%%%%%%%%%%
% LCD DISPLAY
%%%%%%%%%%%%%%%%%%%%%%%%%%%%%%%%%%%%%%%%%%%%%%%%%%
\newpage
\section{LCD Display}

The scientific calculator uses a 16×2 LCD display as the primary output device. It displays the mathematical expressions entered through the keypad as well as the calculated results. The LCD is interfaced with the Arduino Uno in 4-bit mode, reducing the number of I/O pins required while maintaining reliable communication.

\begin{figure}[H]
    \centering
    % Replace lcd_display.jpg with your image filename
    \includegraphics[width=0.65\textwidth]{display.jpeg}
    \caption{16×2 LCD Display Showing Calculator Output}
\end{figure}

%%%%%%%%%%%%%%%%%%%%%%%%%%%%%%%%%%%%%%%%%%%%%%%%%%
% FEATURES
%%%%%%%%%%%%%%%%%%%%%%%%%%%%%%%%%%%%%%%%%%%%%%%%%%

\section{Features}

The developed scientific calculator provides the following features:

\begin{itemize}

\item Basic arithmetic operations.

\item Scientific functions such as sine, cosine and tangent.

\item Exponential and logarithmic calculations.

\item LCD-based output display.

\item Matrix keypad for user input.

\item Breadboard-based modular design.

\item Easy debugging and hardware modification.

\end{itemize}

%%%%%%%%%%%%%%%%%%%%%%%%%%%%%%%%%%%%%%%%%%%%%%%%%%
% APPLICATIONS
%%%%%%%%%%%%%%%%%%%%%%%%%%%%%%%%%%%%%%%%%%%%%%%%%%
\newpage
\section{Applications}

The developed calculator can be used in the following applications:

\begin{itemize}

\item Educational laboratories.

\item Embedded systems learning.

\item Engineering mathematics demonstrations.

\item Microcontroller-based projects.

\item Student mini-projects.

\end{itemize}

%%%%%%%%%%%%%%%%%%%%%%%%%%%%%%%%%%%%%%%%%%%%%%%%%%
% ADVANTAGES
%%%%%%%%%%%%%%%%%%%%%%%%%%%%%%%%%%%%%%%%%%%%%%%%%%

\section{Advantages}

\begin{itemize}

\item Simple and cost-effective implementation.

\item Easy to assemble and troubleshoot.

\item Supports multiple mathematical operations.

\item User-friendly LCD interface.

\item Modular breadboard design.

\item Easily expandable for additional functions.

\end{itemize}

%%%%%%%%%%%%%%%%%%%%%%%%%%%%%%%%%%%%%%%%%%%%%%%%%%
% CONCLUSION
%%%%%%%%%%%%%%%%%%%%%%%%%%%%%%%%%%%%%%%%%%%%%%%%%%

\section{Conclusion}

The scientific calculator project successfully demonstrated the implementation of a microcontroller-based calculator using an Arduino Uno, a 16×2 LCD display, and a 5×5 push-button matrix.

The project provided practical experience in Arduino programming, keypad interfacing, LCD communication, and breadboard circuit design. The completed prototype was capable of accepting user inputs, performing mathematical calculations, and displaying the results accurately on the LCD.

Overall, the project strengthened the understanding of embedded systems, digital electronics, and hardware interfacing concepts while providing hands-on experience in building a functional calculator prototype.

\end{document}
